% Volume VII: Deep Learning Mathematics
% Continuity note for Chapter 39.
% Previous dependencies: Chapters 20-23 introduced objectives, gradients, stochastic optimization, and automatic differentiation; Chapters 37-38 introduced neural networks as parameterized function classes and backpropagation.
% New durable concepts: nonconvex loss landscapes, interpolation, Hessian directions, initialization and variance propagation, normalization, residual gradient paths, stochastic gradient noise, learning-rate schedules, implicit bias, minimum-norm selection, margin bias, and double-descent caveats.
% Forward hooks: Chapter 40 uses these training ideas inside modern architectures; Chapter 41 reuses nonconvex games and likelihood objectives for generative models; Chapter 42 uses optimization geometry to explain robustness and interpretability limits.
\section{Chapter 39: Deep Optimization and Training Dynamics}
\label{ch:deep-optimization-training-dynamics}

Deep learning does not merely ask whether a function class contains a good predictor. It asks whether a practical algorithm can find a useful parameter vector inside a huge, nonconvex, often overparameterized space. The same architecture can train smoothly or fail completely depending on initialization, normalization, learning-rate schedule, gradient noise, and the implicit preferences of the optimizer.

\paragraph{Chapter problem.}
How do gradient-based algorithms move through deep-network loss landscapes, and why do initialization, normalization, residual paths, stochasticity, and implicit bias shape the solutions they find?

\paragraph{Section map.}
Section 39.1 studies loss landscapes through gradients, Hessians, saddle points, flat directions, interpolation, and training trajectories. Section 39.2 derives signal-preserving initialization rules and explains symmetry breaking. Section 39.3 explains normalization and residual architectures as mechanisms that improve gradient flow and conditioning. Section 39.4 studies stochastic training dynamics, learning-rate schedules, implicit bias, and generalization caveats such as double descent.

\subsection{39.1 Loss landscapes}

\paragraph{Guiding question.}
When a deep network has millions or billions of parameters, what does it mean to optimize its loss, and why is convex intuition only a partial guide?

\paragraph{Beginner-facing intuition.}
A supervised deep network is trained by changing parameters until predictions improve on a dataset. If the parameter vector is \(\theta\in\mathbb{R}^p\), a typical empirical loss is
\[
L(\theta)=\frac{1}{n}\sum_{i=1}^n \ell(f_\theta(x_i),y_i),
\]
where \(f_\theta\) is the network, \((x_i,y_i)\) are examples, and \(\ell\) is a loss such as squared error or cross-entropy. The loss landscape is the graph of \(L\) over parameter space.

For a convex quadratic, the landscape is one bowl and gradient descent heads toward its bottom. Deep-network losses are usually nonconvex because the parameters enter through many nested nonlinear maps. The landscape may contain saddle points, nearly flat valleys, symmetries, and many parameter vectors that compute almost the same function. The goal is not to find a mythical unique global bottom. The practical goal is to follow a stable training trajectory to a parameter vector whose predictions generalize.

\paragraph{Formal definitions and notation.}
Let \(L:\mathbb{R}^p\to\mathbb{R}\) be differentiable. Its gradient is
\[
\nabla L(\theta)=
\begin{bmatrix}
\partial L/\partial \theta_1\\
\vdots\\
\partial L/\partial \theta_p
\end{bmatrix},
\]
and, when second derivatives exist, its Hessian is
\[
H(\theta)=\nabla^2 L(\theta)\in\mathbb{R}^{p\times p}.
\]
A point \(\theta^\star\) is a \emph{critical point} if \(\nabla L(\theta^\star)=0\). It is a \emph{local minimum} if \(L(\theta^\star)\leq L(\theta)\) for all \(\theta\) in some neighborhood. It is a \emph{saddle point} if it is critical but nearby directions include both non-increasing and non-decreasing behavior. A model \emph{interpolates} a training set when it achieves zero or nearly zero training loss.

\paragraph{Core mechanism: the Hessian gives the local landscape.}
For a twice differentiable loss, Taylor approximation near \(\theta\) gives
\[
L(\theta+u)
=L(\theta)+\nabla L(\theta)^\top u+\frac12 u^\top H(\theta)u+o(\|u\|^2).
\]
At a critical point, the linear term vanishes. If all Hessian eigenvalues are positive, every small direction increases the loss to second order, so the point is a strict local minimum. If all are negative, it is a local maximum. If some are positive and some are negative, the point is a saddle. If many eigenvalues are near zero, the landscape is locally flat in many directions; this is common in overparameterized networks because many parameter perturbations change the represented function only weakly.

\paragraph{Concrete example.}
Consider the two-parameter model
\[
\widehat y=abx
\]
with one training example \(x=1,y=1\) and squared loss
\[
L(a,b)=(ab-1)^2.
\]
This loss is not convex in \((a,b)\). Every pair satisfying \(ab=1\) is a global minimizer, so there is an entire valley of zero training loss:
\[
(1,1),\quad (2,1/2),\quad (-1,-1),\quad \ldots
\]
At \((0,0)\), the gradient is zero because
\[
\frac{\partial L}{\partial a}=2(ab-1)b,\qquad
\frac{\partial L}{\partial b}=2(ab-1)a.
\]
But \((0,0)\) is not a minimum: moving to \(a=b=\varepsilon\) gives \(L(\varepsilon,\varepsilon)=(\varepsilon^2-1)^2<1=L(0,0)\) for small nonzero \(\varepsilon\). This tiny example already shows nonconvexity, a saddle-like critical point, and many equivalent solutions.

\paragraph{Computational neuroscience example.}
Suppose an RNN is fit to neural population trajectories during a reaching task. The parameters include recurrent weights, input weights, and readout weights. A low loss may be achieved by many different recurrent circuits that generate similar output trajectories. Some parameter directions change the internal dynamics without changing the decoded reach, so the Hessian can have flat directions. A fitted RNN is therefore a model of possible dynamics, not a unique recovered biological circuit.

\paragraph{AI/ML example.}
In image classification, a modern network can drive training cross-entropy nearly to zero. Training curves usually show loss decreasing along a path, not a direct jump to the best possible model. Saddle regions and plateaus slow progress; sharp directions can make large learning rates unstable; flat valleys can allow many parameter values with similar training loss. The optimizer, batch order, and schedule influence which valley is reached.

\paragraph{Common misconceptions and failure modes.}
One misconception is that nonconvex means hopeless. In deep learning, overparameterization often creates many low-loss routes. Another is that a lower training loss always means a better model; interpolation can coexist with poor test performance if the data, labels, or inductive bias are unfavorable. A third is to classify all local minima as bad traps. In high-dimensional neural networks, saddle points and ill-conditioned flat regions are often more important training obstacles than isolated bad minima. Finally, Hessian curvature is local; it does not describe the whole landscape.

\paragraph{Exercises.}
\begin{enumerate}[label=\textbf{39.1.\arabic*.}, leftmargin=*]
\item \textbf{Mechanical.} For \(L(a,b)=(ab-1)^2\), compute \(\nabla L(a,b)\) and verify that every point with \(ab=1\) has zero loss.
\item \textbf{Proof/derivation.} Use the second-order Taylor expansion to show why a critical point with a Hessian eigenvector \(v\) satisfying \(v^\top H v<0\) cannot be a strict local minimum.
\item \textbf{Numerical/computational.} Evaluate \(L(a,b)=(ab-1)^2\) at \((0,0)\), \((0.1,0.1)\), \((1,1)\), and \((2,0.5)\). Which points interpolate the one-example dataset?
\item \textbf{Interpretation.} In a fitted RNN for neural trajectories, why might two very different recurrent matrices have almost the same training loss?
\item \textbf{Misconception/debugging.} A student says, ``Since the loss is nonconvex, gradient descent cannot work.'' Give a precise correction using overparameterization, local curvature, and empirical training trajectories.
\end{enumerate}

\subsection{39.2 Initialization}

\paragraph{Guiding question.}
How should neural-network weights be chosen before training so that units are not symmetric copies and signals neither vanish nor explode?

\paragraph{Beginner-facing intuition.}
Initialization is the starting point of optimization. If every hidden unit starts with identical weights, then backpropagation gives identical updates to those units, so they remain identical and fail to learn different features. Randomness breaks this symmetry. But random weights that are too small erase signals as depth grows, while random weights that are too large amplify signals and gradients uncontrollably.

Good initialization is therefore a scale-matching problem. Each layer should transform a typical input into an output with comparable variance, and its backward pass should transmit gradients at a comparable scale. Xavier and He initializations are simple versions of this idea.

\paragraph{Formal definitions and notation.}
Consider one layer
\[
z_i=\sum_{j=1}^{d_{\mathrm{in}}}W_{ij}h_j,\qquad
h_i^+=\phi(z_i),
\]
where \(h\in\mathbb{R}^{d_{\mathrm{in}}}\), \(W\in\mathbb{R}^{d_{\mathrm{out}}\times d_{\mathrm{in}}}\), and \(\phi\) is an activation. Assume the entries \(W_{ij}\) are independent with mean \(0\) and variance \(\sigma_W^2\), and that the inputs \(h_j\) are approximately independent with mean \(0\) and variance \(q\). The \emph{fan-in} is \(d_{\mathrm{in}}\), and the \emph{fan-out} is \(d_{\mathrm{out}}\).

\paragraph{Core derivation: variance propagation.}
Under the independence assumptions,
\[
\operatorname{Var}(z_i)
=\operatorname{Var}\left(\sum_{j=1}^{d_{\mathrm{in}}}W_{ij}h_j\right)
=\sum_{j=1}^{d_{\mathrm{in}}}\operatorname{Var}(W_{ij}h_j)
=d_{\mathrm{in}}\sigma_W^2 q.
\]
For a roughly linear or tanh-like layer near zero, preserving forward variance suggests
\[
d_{\mathrm{in}}\sigma_W^2\approx 1.
\]
Balancing forward and backward variance leads to the Xavier scale
\[
\sigma_W^2\approx \frac{2}{d_{\mathrm{in}}+d_{\mathrm{out}}}.
\]
For ReLU, if \(z\) is roughly symmetric around zero, then \(\phi(z)=\max(0,z)\) keeps the positive half of the activity. The clean half-factor is for the second moment:
\[
\mathbb{E}[\phi(z)^2]=\frac12\mathbb{E}[z^2]
\]
when the distribution of \(z\) is symmetric around zero. For \(z\sim\mathcal{N}(0,s^2)\), the centered variance is instead
\[
\operatorname{Var}(\phi(z))=s^2\left(\frac12-\frac{1}{2\pi}\right),
\]
because ReLU outputs have positive mean. He initialization is therefore best understood as preserving the uncentered second moment that enters the next random signed affine layer. Setting
\[
\frac12 d_{\mathrm{in}}\sigma_W^2 q\approx q
\]
suggests the He scale
\[
\boxed{\sigma_W^2\approx \frac{2}{d_{\mathrm{in}}}.}
\]
These formulas are approximations, not laws of nature; they depend on activation, architecture, normalization, and distributional assumptions.

\paragraph{Concrete example.}
Suppose a ReLU layer has \(d_{\mathrm{in}}=100\). He initialization chooses
\[
\sigma_W^2=\frac{2}{100}=0.02,\qquad \sigma_W\approx 0.141.
\]
If the input coordinates have variance \(q=1\), then the preactivation variance is approximately
\[
d_{\mathrm{in}}\sigma_W^2 q=100(0.02)(1)=2.
\]
After ReLU, the second moment \(\mathbb{E}[\phi(z)^2]\) is roughly \(1\). The centered variance is smaller because the ReLU activation has positive mean, but the next random signed affine layer receives activations of comparable second-moment scale.

\paragraph{Computational neuroscience example.}
Random recurrent networks used as models of cortical dynamics often scale recurrent weights by the number of neurons, for example with variance proportional to \(1/N\). This keeps total recurrent input from growing uncontrollably as the population size \(N\) increases. The spectral radius of the recurrent matrix then influences whether activity decays, persists, or becomes unstable. Such initialization choices are modeling assumptions about a dynamical regime, not direct measurements of synapses.

\paragraph{AI/ML example.}
In a deep ReLU classifier without normalization, initializing all weights with variance \(1\) would usually make activations explode across layers. Initializing with extremely tiny variance would make activations and gradients vanish. He initialization gives a practical starting scale. In transformers, initialization interacts with residual connections and layer normalization, so the useful scale is architecture-dependent.

\paragraph{Common misconceptions and failure modes.}
Zero initialization is not merely ``slow''; in hidden layers it preserves permutation symmetry among units. Variance formulas are not exact because real activations are dependent, biased, and changed by training. A good initialization does not guarantee good optimization; it only avoids immediate signal-scale failure. Finally, initializing recurrent weights near a desired spectral radius does not make the network biologically realistic by itself.

\paragraph{Exercises.}
\begin{enumerate}[label=\textbf{39.2.\arabic*.}, leftmargin=*]
\item \textbf{Mechanical.} For a ReLU layer with fan-in \(256\), compute the He initialization variance and standard deviation.
\item \textbf{Proof/derivation.} Under the independence assumptions in this section, derive \(\operatorname{Var}(z_i)=d_{\mathrm{in}}\sigma_W^2q\).
\item \textbf{Numerical/computational.} A three-layer linear network multiplies variance by \(1.5\) at each layer. If the input variance is \(2\), what is the output variance after three layers?
\item \textbf{Interpretation.} Explain why identical initialization of two hidden units makes them receive identical gradients in a fully connected layer with the same activation.
\item \textbf{Misconception/debugging.} A student initializes a \(50\)-layer ReLU network with weights of standard deviation \(1\) and observes numerical overflow. Identify the scale mistake and propose a variance-preserving alternative.
\end{enumerate}

\subsection{39.3 Normalization and residual architectures}

\paragraph{Guiding question.}
How do normalization layers and residual connections make deep networks easier to train without changing the basic goal of minimizing a loss?

\paragraph{Beginner-facing intuition.}
Deep networks are compositions. If each layer changes the scale or conditioning of the signal unpredictably, the next layer receives a moving target. Normalization layers recenter and rescale activations so later computations see more controlled inputs. Residual connections add an identity path, so a block can learn a correction rather than an entirely new transformation.

The practical benefit is gradient flow. In a plain deep stack, the backward signal is repeatedly multiplied by layer Jacobians. If many singular values are less than one, gradients vanish; if many are greater than one, gradients explode. Residual paths provide a direct route through the identity, and normalization reduces scale drift.

\paragraph{Formal definitions and notation.}
For a vector \(a=(a_1,\ldots,a_m)\), a basic normalization operation is
\[
\mu=\frac1m\sum_{j=1}^m a_j,\qquad
\sigma^2=\frac1m\sum_{j=1}^m(a_j-\mu)^2,
\]
\[
\widehat a_j=\frac{a_j-\mu}{\sqrt{\sigma^2+\epsilon}},
\qquad
y_j=\gamma_j\widehat a_j+\beta_j.
\]
The small \(\epsilon>0\) prevents division by zero, and \(\gamma,\beta\) are learned scale and shift parameters.

Batch normalization computes statistics over a minibatch and feature dimension chosen by the architecture; it uses running statistics at inference time. Layer normalization computes statistics within each example across features and is common in transformers.

A residual block has the form
\[
\boxed{h_{\ell+1}=h_\ell+F_\ell(h_\ell;\theta_\ell),}
\]
where \(h_\ell\in\mathbb{R}^d\) and \(F_\ell:\mathbb{R}^d\to\mathbb{R}^d\) is a learned transformation. If dimensions differ, a projection may replace the identity path.

\paragraph{Core derivation: the residual gradient path.}
Let \(g_{\ell+1}=\nabla_{h_{\ell+1}}L\). Since
\[
h_{\ell+1}=h_\ell+F_\ell(h_\ell),
\]
the Jacobian with respect to \(h_\ell\) is
\[
\frac{\partial h_{\ell+1}}{\partial h_\ell}
=I+J_{F_\ell}(h_\ell).
\]
Therefore the backward gradient is
\[
\boxed{
g_\ell
=\left(I+J_{F_\ell}(h_\ell)\right)^\top g_{\ell+1}.
}
\]
The identity term means that part of the gradient can pass directly backward even if the learned block has small derivatives. This does not remove all optimization problems, but it changes the product of Jacobians from a product of arbitrary transformations into a product of identity-plus-correction transformations.

\paragraph{Concrete example.}
Let \(a=(2,4,6)\). Then
\[
\mu=4,\qquad
\sigma^2=\frac{(2-4)^2+(4-4)^2+(6-4)^2}{3}=\frac83.
\]
Ignoring \(\epsilon\), the normalized values are approximately
\[
\widehat a=(-1.225,0,1.225).
\]
If \(\gamma=(2,2,2)\) and \(\beta=(1,1,1)\), then
\[
y\approx(-1.449,1,3.449).
\]
Normalization fixes mean and variance before the learned affine scale and shift restore representational flexibility.

\paragraph{Computational neuroscience example.}
Divisive normalization in sensory neuroscience describes response gain controlled by activity in a pool of neurons. This resembles the scale-control intuition behind normalization layers, but the mathematical objects differ: batch normalization is an engineering operation over examples and features, while biological normalization is a circuit-level response model. The analogy is useful only when the axes and mechanisms are stated.

\paragraph{AI/ML example.}
Residual connections made very deep CNNs trainable by letting blocks refine feature maps. Layer normalization is central in transformers because sequence length and batch composition vary, and normalization within each token representation gives stable feature scales. In both cases, the architecture changes the conditioning of the optimization problem, not the definition of the supervised objective.

\paragraph{Common misconceptions and failure modes.}
Normalization does not make a network invariant to all input scaling; it normalizes selected axes at selected locations. Batch normalization has different behavior during training and inference, so stale running statistics can cause deployment errors. Residual connections do not guarantee stability if \(F_\ell\) is too large or the learning rate is poorly chosen. A residual block is not automatically biologically realistic merely because it resembles a recurrent correction.

\paragraph{Exercises.}
\begin{enumerate}[label=\textbf{39.3.\arabic*.}, leftmargin=*]
\item \textbf{Mechanical.} Normalize \(a=(1,1,3,5)\) using the formula in this section with \(\epsilon=0\), \(\gamma_j=1\), and \(\beta_j=0\).
\item \textbf{Proof/derivation.} For \(h_{\ell+1}=h_\ell+F_\ell(h_\ell)\), derive \(g_\ell=(I+J_{F_\ell})^\top g_{\ell+1}\).
\item \textbf{Numerical/computational.} Suppose a plain layer multiplies a one-dimensional gradient by \(0.8\). After \(20\) layers, what is the multiplier? If each residual layer has multiplier \(1+0.02\), what is the multiplier after \(20\) layers?
\item \textbf{Interpretation.} Explain why batch normalization can behave differently for a minibatch of size \(4\) than for a minibatch of size \(256\).
\item \textbf{Misconception/debugging.} A model trains well but fails at inference after batch normalization is added. Name one possible statistics-related cause and how to check it.
\end{enumerate}

\subsection{39.4 Implicit bias of optimization}

\paragraph{Guiding question.}
When many parameter vectors fit the training data, which one does gradient-based training tend to choose, and how do stochasticity and schedules affect that choice?

\paragraph{Beginner-facing intuition.}
In overparameterized learning, the training equations may have many exact solutions. The optimizer then does more than minimize the explicit loss. Its update rule, initialization, batch noise, learning-rate schedule, and parameterization create an \emph{implicit bias}: a preference among solutions that is not written as an explicit penalty in the objective.

This matters for generalization. Two networks with the same training loss can behave differently on new data. Some differences come from explicit choices such as weight decay or data augmentation. Others come from the training algorithm itself.

\paragraph{Formal definitions and notation.}
Full-batch gradient descent uses
\[
\theta_{t+1}=\theta_t-\eta_t\nabla L(\theta_t),
\]
where \(\eta_t>0\) is the learning rate. Stochastic gradient descent uses a minibatch \(B_t\):
\[
\theta_{t+1}=\theta_t-\eta_t g_{B_t}(\theta_t),
\qquad
g_{B_t}(\theta)=\frac{1}{|B_t|}\sum_{i\in B_t}\nabla_\theta \ell_i(\theta).
\]
If the minibatch is sampled uniformly, then
\[
\mathbb{E}[g_{B_t}(\theta)]=\nabla L(\theta),
\]
and the gradient noise is
\[
\xi_t=g_{B_t}(\theta_t)-\nabla L(\theta_t).
\]
A learning-rate schedule is the sequence \((\eta_t)\), often including warmup, constant phases, step decay, cosine decay, or exponential decay.

\paragraph{Core derivation: minimum-norm bias in underdetermined linear regression.}
Let \(X\in\mathbb{R}^{n\times p}\) with \(p>n\), and suppose the linear system
\[
X\theta=y
\]
has many solutions. Consider squared loss
\[
L(\theta)=\frac12\|X\theta-y\|_2^2
\]
and gradient descent initialized at \(\theta_0=0\). The gradient is
\[
\nabla L(\theta)=X^\top(X\theta-y),
\]
which always lies in the row space of \(X\). Since \(\theta_0=0\) lies in the row space plus no null-space component, every iterate remains in the row space of \(X\). Among all exact solutions, exactly one lies in the row space:
\[
\boxed{\theta^\star=X^\top(XX^\top)^{-1}y}
\]
when \(XX^\top\) is invertible. This is the minimum Euclidean norm interpolating solution. Thus gradient descent has selected a solution even though the objective did not explicitly say ``minimize norm.''

For separable logistic regression, a related but different result holds: under suitable conditions, gradient descent drives the norm to infinity while the direction approaches the maximum-margin separator. This is another implicit bias, now toward margin rather than finite norm.

\paragraph{Concrete example.}
Let
\[
X=\begin{bmatrix}1&1\end{bmatrix},\qquad y=1.
\]
The equation \(X\theta=y\) is
\[
\theta_1+\theta_2=1,
\]
which has infinitely many solutions. The minimum-norm solution is \((1/2,1/2)\). Starting from \((0,0)\), the gradient is always proportional to \(X^\top=(1,1)^\top\), so gradient descent moves along the line \(\theta_1=\theta_2\) and converges to \((1/2,1/2)\), not to \((1,0)\) or \((0,1)\).

\paragraph{Computational neuroscience example.}
A neural decoder may map a high-dimensional population response to a movement variable. If many weight vectors decode the training trials equally well, gradient descent from small initialization can prefer a small-norm decoder that spreads weight across correlated neurons. This is an algorithmic preference, not proof that the biological system uses the same distributed readout.

\paragraph{AI/ML example.}
Large classifiers often interpolate their training data. Generalization then depends on data distribution, architecture, augmentation, optimizer, weight decay, training time, and implicit bias. The \emph{double-descent} phenomenon describes test error that can decrease, rise near an interpolation threshold, and decrease again as model size grows. It is not a universal law; it is an empirical and theoretical pattern that appears under some data, model, and training regimes.

\paragraph{Common misconceptions and failure modes.}
Implicit bias is not the same as explicit regularization. Weight decay adds a term or update effect by design; implicit bias comes from the algorithm and parameterization. SGD noise is not always helpful noise; too much noise or too large a learning rate can prevent convergence. A decaying learning rate can help settle into a basin, but decaying too early can freeze a poor solution. Double descent does not mean ``bigger is always better''; compute, data quality, distribution shift, and evaluation metric matter.

\paragraph{Exercises.}
\begin{enumerate}[label=\textbf{39.4.\arabic*.}, leftmargin=*]
\item \textbf{Mechanical.} For \(X=[1\ 1]\) and \(y=1\), list three exact solutions to \(X\theta=y\). Which has minimum Euclidean norm?
\item \textbf{Proof/derivation.} Show that, for squared linear regression, \(\nabla L(\theta)=X^\top(X\theta-y)\) lies in the row space of \(X\).
\item \textbf{Numerical/computational.} Starting from \(\theta_0=(0,0)\), learning rate \(\eta=0.25\), and loss \(L(\theta)=\frac12(\theta_1+\theta_2-1)^2\), compute two gradient-descent steps.
\item \textbf{Interpretation.} Explain why two interpolating neural networks can have identical training loss but different test behavior.
\item \textbf{Misconception/debugging.} A student observes double descent in one experiment and concludes that adding parameters always improves generalization. State at least three conditions or caveats that make this conclusion invalid.
\end{enumerate}

\paragraph{Closing bridge.}
This chapter treated training as a dynamical process through parameter space. Chapter 40 turns to the architectures that define those parameter spaces: convolutional, recurrent, attention-based, and graph-based networks. The same optimization principles will reappear there, but each architecture adds its own tensor shapes, invariances, memory structure, and failure modes.

\clearpage
